\documentclass[english,twoside, openright, hidelinks, 11 pt, class=report,crop=false]{standalone}
\input{../../preamb}
\input{../bm}

\usepackage{xr}
\externaldocument{tri_opg_bm}


\begin{document}
\subsection*{Kapittel \ref{Trigonometri}}	
\footnotesize
For alle svar tas det for gitt at $ n\in \mathbb{Z} $.

\textsl{Merk}: Uttrykkene for løsninger av trigonometriske ligninger kan se forskjellige ut, men gi de samme verdiene av $ x $. For eksempel vil $ {x= 2\pi n-\frac{\pi}{4}} $ være den samme løsningen som $ {x = \frac{7\pi}{4}+2\pi n} $ fordi $ {-\frac{\pi}{4}+2\pi = \frac{7\pi}{4}}$. Vi kan alltid trekke ut heltallsfaktorer av $ n $-leddet for å endre på uttrykk, for å sjekke om ditt svar er riktig bør du derfor først sjekke at ditt $ n $-ledd er i overensstemmelse med fasit.

\opr{rads} 
\abch{
\item $ \frac{\pi}{3} $
\item $ \frac{\pi}{12} $
\item $\frac{4\pi}{9}$
\item $\frac{\pi}{10}$
}

\opr{grad}
\abch{
\item $ 165^\circ $ 
\item $ 330^\circ $
\item $108^\circ$
\item $90^\circ$
}

\opr{r2v25d1opg7a}
\abch{
\item $\frac{4}{3}$
\item $\frac{240^\circ}{\pi}$
}


\opr{pytg} \selos

\opr{tanx} 
\abch{
\item $ 0  $ 
\item $- \frac{1}{\sqrt{3}} $
}

\opr{r2h24d1opg4}
$\cos x=-\frac{\sqrt{15}}{4}$, $\tan x=\frac{\sqrt{15}}{15}$

\opr{r2v25d1opg7b}
$\cos x = \frac{1}{\sqrt{5}}$, $\sin x=\frac{2}{\sqrt{5}}$

\opr{r2v26d1opg3}
$\sin v =- \frac{\sqrt{5}}{3}$, $\tan v=-\frac{\sqrt{5}}{2}$

\opr{r2v26d1opg4}
$f(x)=-2\cos\left(2x\right)-1$

\opr{kvadrant}\\
\begin{tabular}{@{}l|c| c| c|c|}

	&1. kvadrant & 2. kvadrant & 3. kvadrant & 4. kvadrant\\ \hline
	$ \sin x $& +& +& $ - $&$ - $\\ \hline
	$ \cos x $& +& $ - $& $ - $&+\\ \hline
	$ \tan x $& +& $ - $& +&$ - $\\\hline
\end{tabular}\vsk

\opr{trigverd}
\abch{
\item $ -\frac{1}{2} $ 
\item $ -1 $ 
\item $ 0 $ 
\item $ -\sqrt{3} $
}


\opr{averdier} 
\abch{
\item 0 
\item $ \frac{\pi}{3} $ 
\item $ \pi $
\item $ \frac{3\pi}{4} $ 
\item $ \frac{\pi}{4} $ 
\item $ \frac{\pi}{6} $
}


\opr{bruk1} \selos

\opr{sin2xopg} 
\abch{
\item Se side \pageref{sin2xbevis}. 
\item Se løsningsforslag.
}


\opr{cossomsin} $ \sin(3x) $

\opr{cossinsin} 
$2\sin\left(2x+\frac{\pi}{6}\right)  $ 

\opr{sincos0} Se løsningsforslag.

\opr{triligns}
\abch{
\item $ x = \pm \frac{\pi}{4}+2\pi n $ 
\item $ x=\frac{3}{2} \;\vee\; x = \frac{9}{2} $ 
\item $ x=\frac{1}{3}(\frac{\pi}{6}+2\pi n)\;\vee\; x= \frac{1}{3}(\frac{5\pi}{6}+2\pi n) $ 
\item $ x=\frac{\pi}{3}(3n+1)\;\vee\; x= \frac{\pi}{6}(6 n+1)$ 
\item $\frac{1}{4}(\pi n-\frac{\pi}{3} )$
\item $x\in\{\frac{1}{2}, \frac{11}{2}, \frac{13}{2}\}$
}


\opr{asinbcos0o}
\abch{
\item $ \frac{\pi}{6}+\pi n $ 
\item $ \pi n-\frac{\pi}{3} $ 
\item $ \frac{1}{2}\left(\pi n-\frac{\pi}{4}\right) $
}

\opr{kombo}
\abch{
\item $ x=2\pi n-\frac{\pi}{4} $
\item $ x\pi^2(4n-1)\;\vee\; x=2\pi\left(\frac{\pi}{6}+2\pi n\right) $
\item $x\in\left\{\frac{\pi}{3}, \frac{4}{3}\pi\right\}$ 
}


\opr{abctri}
\abch{
\item ${ x= 2\pi n-\frac{\pi}{2} \;\vee\; x = \frac{\pi}{6}+2\pi n \;\vee\; x = \frac{5\pi}{6}+2\pi n}$ 
\item ${ x=\frac{1}{3}(\pi +2\pi n) }$ 
\item ${ x = \pm \frac{3\pi}{4}+2\pi n  }$ 
\item ${ x=\frac{1}{3}+n }$ 
\item ${ x=\frac{1}{3}(\pi +2\pi n) }$
}


\opr{kvadopg}
\abch{
\item $ x = \pm \frac{\pi}{6}+\pi n $ 
\item $ \pm \pi + 4\pi n$ (eventuelt $ x=\pi +2\pi n) $)
}

\opr{r2v23d1opg3}
$x=\frac{\pi}{3}$, $x=\frac{2\pi}{3}$, $x=\frac{4\pi}{3}$, $x=\frac{5\pi}{3}$


\opr{cosmaks} \selos

\opr{cosfunko} 
\abch{
\item $ P = \frac{2\pi}{3} $ 
\item $ f_{maks}=7 $, $ f_{min}=1 $ 
\item $ f $ har minimum for $ x= \frac{1}{3}\left(2\pi n-\frac{\pi}{12}\right) $ og maksimum for $ x=\frac{1}{3}(2\pi n-\frac{11\pi}{12}) $
}

\opr{opgtrifinnmaks}
$x=\frac{2n-1}{3}\pi$

\opr{sinmaks}
\selos

\opr{sinfnpkt} 
\abch{
\item $ P=4 $ 
\item $ (-1, 3) $ og $ (3, 3) $ 
\item $x=-\frac{11}{3}$, $ x=-\frac{7}{3}, x=\frac{5}{3} $ og $ x=\frac{1}{3} $
}


\opr{skissin} \\
\includegraphics[scale=0.5]{\figp{skissin}}

\opr{opgtriomskriv}
\abch{
\item \selos
\item \selos
}

\opr{opgtriomskriv1}
$f(x)=2 \sin \left(3x-\frac{\pi}{2}\right)+1 $

\opr{sincosfig}
\abch{
	\item $ 3\cos(\pi x - \pi) +2 $ 
	\item $ 3\sin\left(\pi x-\frac{\pi}{2}\right) +2 $
}

\opr{tanfopg}
\selos

\opr{r2v24d1opg5}
\abch{
\item $x\in\{3, 7, 15, 19\}$
\item amplitude 2, likevektslinje $y=-1$, periode 12, forskyvning 2
}

\opr{r2v25d1opg4}
\abch{
\item Amplitude $2\sqrt{3}$, likevektslinje $d=0$, periode $\pi$ og faseforskyvning $-\frac{\pi}{12}$
\item $x=\frac{\pi}{3}$
\item $x\in\left\{\frac{\pi}{3}, \pi, \frac{4\pi}{3}\right\}$
}

\opr{opgtriderfunk}

\grubr{udir}
\selos

\grubr{opgtricos2}
\selos

\grubr{r2h24d1opg5}
\abch{
\item $f(x)=2\cos\left(\frac{\pi}{4}x-\pi\right)-1$.
\item $x\in\{\frac{8}{3}, \frac{16}{3}\}$. Vi fant nullpunktene til funksjonen.
}



\grubr{eksmpr2v23d2opg2}
\abch{
\item $x\in\left[0, \frac{\pi}{6}\right)\cup\left(\frac{5\pi}{6},\pi\right]$
\item $b\in\left[0, \frac{\sqrt{3}}{4}\right)\cup\left(\frac{3\sqrt{3}}{4}, \infty\right)$
}

\grubr{r2h25d1opg4}
To løsninger for $a=\frac{1}{2}$, tre løsninger for $a=\pm 1$. Fire løsninger for alle andre $a\in[-1,1]$. Ingen løsninger for $|a|>1$.

\grubr{r2v23d1opg5}
\abch{
\item \selos
\item \selos
\item \selos
}

\grubr{opgtangr} \selos


\end{document}

